\section{Welfare and sale terms}\label{sec:11_welfare}

This chapter holds the incumbent fixed and asks who gains when the challenger sees the price. It then changes the
sale terms: the payment rule and the reserve. Write \(\mathcal R\) for expected target proceeds and \(\mathcal W\) for
expected acquisition surplus net of preparation costs. Acquisition surplus excludes transfers between bidders,
shareholders and traders. \paper{Sec.~7; Appendix~A.7}

\subsection{Access to prices: two equilibria with the same orders}\label{sec:11_welfare-access}
\paper{Sec.~7.1; Prop.~A.8}

\begin{intuition}[Same orders, better preparation decisions]
At the strong incumbent the investor trades fully even when the challenger cannot see the price. So access leaves
the orders and the trading costs unchanged. It changes one decision: a high-cost challenger now prepares after a
purchase, where its expected gross profit exceeds its cost.
\end{intuition}

Hold the incumbent at \(F_1\) under the hypotheses of Proposition~2, and solve two economies. With access, the
equilibrium is the strong tuple \(\mathcal E_1\) of \Cref{sec:06_main_result-strong}: full orders, posteriors
\((m,\pi,M)\) and entry \((\rho,\rho,1)\). With a hidden price, the challenger's belief stays at \(\pi\), and
\(B_1(\pi)\le B_0(\pi)<c_H\), so entry is \(\rho\) after every order. The trading bound holds for every entry in
\([\rho,1]\) (\Cref{sec:05_investor-bounds}), so full orders are again unique; \Cref{sec:07_controls-hidden} solves
this economy. The orders, and hence the trading costs, coincide at every realization.

\begin{resultbox}[Access to prices \paper{Prop.~A.8}]
Under the hypotheses of Proposition~2, at \(F_1\), with superscript \(1\) for access, \(0\) for the hidden price and
\(w_j=t_j-t_0\):
\begin{align*}
\mathcal R^1-\mathcal R^0&=(1-\rho)\bigl\{\pi(b+\lambda)w_H+(1-\pi)\,b\,w_L\bigr\}>0,\\
\mathcal W^1-\mathcal W^0&=(1-\rho)(b+\pi\lambda)\bigl[B_1(M)-c_H+pF_1(p)\bigr]>0.
\end{align*}
\end{resultbox}

\paragraph{The proceeds gain.} Couple the two economies on the class, the value, the cost, the slot identity and the
liquidity order. Only high-cost preparation after a purchase differs. A purchase has probability \(b+\lambda\) in class
\(H\) and \(b\) in class \(L\) (\Cref{sec:06_main_result-entry}). The cost is high with probability \(1-\rho\). Each added
preparation raises expected proceeds from \(t_0\) to \(t_j\), that is by \(w_j\). Multiply the three factors and weight
the classes by \(\pi\) and \(1-\pi\). Both terms are nonnegative, and \(w_H=w_L+\Delta_1\ge\Delta_1>0\), so the gain is positive.

\paragraph{The surplus gain.} Without a challenger the allocation value is \(\mathcal V_0(R)=R\,\1\{R\ge p\}\). Truthful
bids give the target to the highest admissible value, so an admissible challenger with value \(\theta\) adds
\[
\max\{\mathcal V_0(R),\theta\}-\mathcal V_0(R)=(\theta-\max\{p,R\})_++p\,\1\{R<p\}.
\]
If \(R<p\), both sides equal \(\theta\), since \((\theta-p)+p=\theta\). If \(R\ge p\), both sides equal \((\theta-R)_+\).
Take the expectation over \(R\) and over the class at the purchase posterior \(M\). The first term gives \(B_1(M)\) and
the second gives \(pF_1(p)\), the value of the sales that the reserve would otherwise block. Subtract \(c_H\), and
multiply by the probability of an added preparation, \((1-\rho)\Pr(X=1)=(1-\rho)(b+\pi\lambda)\). The bracket is
positive because \(B_1(M)>c_H\) by (H1). The gain is an expectation at the purchase posterior, not a gain at each
realized value.

\paragraph{Benchmark check.} At \(r_1\), \(w_H=9/8\), \(w_L=11/24\) (\Cref{sec:03_auction-terms}) and
\(pF_1(p)=\tfrac12\cdot\tfrac{1/2}{3}=\tfrac1{12}\). With \(\pi=1/2\) and \(\rho=1/4\),
\[
\mathcal R^1-\mathcal R^0=\frac38\Bigl[\frac98(b+\lambda)+\frac{11}{24}b\Bigr]=\reg{p2_access_proceeds_gain},\qquad
\mathcal W^1-\mathcal W^0=\frac34\Bigl(b+\frac\lambda2\Bigr)\Bigl[B_1(M)-6+\frac1{12}\Bigr]=\reg{p2_access_surplus_gain}.
\]
In levels, proceeds rise from \(\reg{p2_base_hidden_proceeds}\) to \(\reg{p2_base_r1_proceeds}\), and net surplus from
\(\reg{p2_base_hidden_surplus}\) to \(\reg{p2_base_r1_surplus}\). The surplus level is
\(\E[\mathcal V_0(R)]=(r^2-p^2)/(2r)\), plus \(g_j+p^2/r\) for each prepared class-\(j\) challenger, minus \(c_H\) for each
high-cost preparation; the low cost is zero. If
the execution cost is a resource cost, both economies lose \(\lambda k\), and the gains do not change. With research
at the cap (\Cref{sec:09_research-strong}), research and trading costs also coincide, and the gains are the same,
\(\reg{res_dR_r1}\) and \(\reg{res_dW_r1}\).

\subsection{Transfers, trade scale and the limits of the gain}\label{sec:11_welfare-transfers}
\paper{Sec.~7.1; Appendix~A.7}

A shareholder who holds through the sale receives \(V_T\), not the interim price, so target shareholders receive the
proceeds gain. The investor's profit is a transfer. In the access economy class \(H\) buys and meets entry one, so
\(K_H=\Delta_1(1-M)\); class \(L\) sells and meets entry \(\rho\), so \(K_L=\rho\Delta_1m\)
(\Cref{sec:05_investor-residual}). Weight by the class probabilities and by the presence probability:
\[
G=\lambda\Delta_1\{\pi(1-M)+(1-\pi)\rho m\}=\reg{p2_access_gross_transfer},\qquad
G/\lambda=\reg{p2_access_gross_conditional}.
\]
Competitive prices give \(\E[V_T-P(X)\mid X]=0\), so \(\E[X(V_T-P(X))]=0\). Split this expectation by who filled the
slot: the informed part is \(G\), so liquidity traders lose \(G\) in expected gross terms. Market makers earn zero, and
the execution cost is separate.

\paragraph{The scale of trade.} Let one order unit be \(z\) shares per share of float, with the same signal
experiment. Without the investor every posterior is \(\pi\) and entry is \(\rho\), so proceeds are \(\mathcal R^0\).
Even if initial shareholders bear the whole liquidity loss, they gain over that economy exactly when
\[
(\mathcal R^1-\mathcal R^0)-zG>0\iff z<\frac{\mathcal R^1-\mathcal R^0}{G}=\reg{p2_access_break_even}.
\]
This holds the signal experiment fixed; it does not solve a game in which the investor chooses its size.

\paragraph{Information, not the price level.} Add a deterministic external dividend of \(\reg{p2_access_proceeds_gain}\)
to the traded claim of the hidden-price economy. \(V_T\) and every price rise by that amount, so mean prices agree
across the economies, while every order payoff, the challenger's information and entry stay the same.

\paragraph{Access outside the recruitment region.} Suppose instead \(B_1(m)<c_H<B_1(\pi)\), with
\(k<\rho\kappa\Delta_1\). Without the price every challenger prepares, so entry is one after every order. With the
price, entry is \((\rho,1,1)\), and orders are full in both economies. Access now \emph{removes} high-cost preparation
after a sale, which has probability \(b\) in class \(H\) and \(b+\lambda\) in class \(L\). The same steps give
\begin{align*}
\mathcal R^1-\mathcal R^0&=-(1-\rho)\{\pi b\,w_H+(1-\pi)(b+\lambda)w_L\}<0,\\
\mathcal W^1-\mathcal W^0&=-(1-\rho)[b+(1-\pi)\lambda]\,[B_1(m)-c_H+pF_1(p)].
\end{align*}
At the strong benchmark with \(c_H=\reg{access_out_cH}\), \(B_1(m)-c_H=\reg{p2_outside_seen_cost_sign_minus}\), but the
bracket is positive because \(pF_1(p)=1/12\). Access changes proceeds by \(\reg{p2_outside_proceeds_gain}\) and surplus
by \(\reg{p2_outside_surplus_gain}\).

\begin{caution}[State the welfare domain]
The positive result holds where the high cost does not pay at the prior. Do not say that price access raises
proceeds or surplus for every cost.
\end{caution}

\subsection{The payment rule: why the sign can change}\label{sec:11_welfare-payment}
\paper{Sec.~7.2; Prop.~A.9}

Proposition~1 belongs to the cash second-price rule (\Cref{sec:03_auction-prop1}). Now change the payment rule. The
institution has zero reserve and verifiable values. After diligence the buyer with the highest value acquires the
target, and a sale to the next buyer at its value is an enforceable fallback. The seller and the winner split the
surplus above the fallback by Nash bargaining, with seller weight \(\omega\in[0,1)\). Keep the point-mass classes, and
let the incumbent's support be \([0,\bar r]\) with \(\ell<\bar r<h\).

\begin{resultbox}[The payment rule \paper{Prop.~A.9}]
\begin{gather*}
T_\omega(R,\theta)=(1-\omega)\min\{R,\theta\}+\omega\max\{R,\theta\},\qquad
g_{\theta,\omega}=(1-\omega)\,\E[(\theta-R)_+],\\
\Delta_\omega=\omega(h-\ell)+(1-2\omega)\,\E[(R-\ell)_+].
\end{gather*}
A stronger incumbent weakly lowers \(g_{\theta,\omega}\), and it moves the spread with the weak sign of \(1-2\omega\);
the move is strict when \(\E[(R-\ell)_+]\) changes strictly and \(\omega\ne1/2\). At \(\omega=1/2\),
\(\Delta_\omega=(h-\ell)/2\) for every incumbent distribution.
\end{resultbox}

\paragraph{The payment.} Let the winner's value be \(V\) and the fallback \(z\le V\). For \(0<\omega<1\) and \(V>z\),
the payment \(P\in[z,V]\) maximizes \((P-z)^\omega(V-P)^{1-\omega}\). Its logarithm has the derivative
\(\omega/(P-z)-(1-\omega)/(V-P)\), which falls strictly in \(P\). Set it to zero:
\[
\omega(V-P)=(1-\omega)(P-z)\iff P=(1-\omega)z+\omega V.
\]
At \(\omega=0\) or \(V=z\) the payment is \(z\), as the formula gives. With a challenger, \(V=\max\{R,\theta\}\) and
\(z=\min\{R,\theta\}\), which is \(T_\omega\). Without one, the fallback is zero and the seller receives \(\omega R\).

\paragraph{Profit and spread.} A challenger wins when \(\theta>R\) and keeps \(\theta-T_\omega=(1-\omega)(\theta-R)\);
this gives \(g_{\theta,\omega}\). For the spread, compare the classes at each \(R<h\). If \(R\le\ell\), both classes win and
pay \((1-\omega)R+\omega\theta\), so the difference is \(\omega(h-\ell)\). If \(\ell<R<h\), the high class wins, and against
the low class the incumbent wins and pays \((1-\omega)\ell+\omega R\):
\[
T_\omega(R,h)-T_\omega(R,\ell)=[(1-\omega)R+\omega h]-[(1-\omega)\ell+\omega R]=\omega(h-\ell)+(1-2\omega)(R-\ell).
\]
Both cases equal \(\omega(h-\ell)+(1-2\omega)(R-\ell)_+\), and the
expectation over \(R\) gives \(\Delta_\omega\). The condition \(\bar r<h\) matters: for \(R\ge h\) both classes lose, and
the difference would be \((1-\omega)(h-\ell)\). Finally, \((\theta-R)_+\) falls and \((R-\ell)_+\) rises in \(R\), so a
first-order strengthening weakly lowers \(g_{\theta,\omega}\) and moves \(\Delta_\omega\) with the sign of \(1-2\omega\).

\begin{intuition}[Sensitivities]
Take \(R>\ell\). The high-class payment \((1-\omega)R+\omega h\) moves with \(R\) at rate \(1-\omega\). Against the low
class the incumbent wins and pays \((1-\omega)\ell+\omega R\), which moves at rate \(\omega\). Competition widens the
gap only if \(1-\omega>\omega\). The zero-reserve second-price auction is \(\omega=0\): all weight on the runner-up's
value.
\end{intuition}

For a uniform incumbent, \(\E[(R-\ell)_+]=(r-\ell)^2/(2r)=\DT(r)\) (\Cref{sec:03_auction-terms}), so
\(\Delta_\omega(r_1)-\Delta_\omega(r_0)=(1-2\omega)(\tfrac23-\tfrac1{60})=\reg{pay_gap_coef}\,(1-2\omega)\). State the scope
before the panel does: the result covers the payment stage only, not the trading and preparation equilibrium under
Nash bargaining.

\subsection{Reserves: preparation, sale and admissible bids}\label{sec:11_welfare-reserves}
\paper{Sec.~7.3; Prop.~A.10; Table~5}

Now let the reserve take any value, so a class value can lie below it. For a point value \(v\) and
\(R\sim U[0,r]\), with \(t_0=p(1-p/r)_+\),
\[
(t_v,g_v,W_v)=\begin{cases}
(t_0,0,0),&v<p,\\
(p,\,v-p,\,1),&r\le p\le v,\\
\bigl(v-\tfrac{v^2-p^2}{2r},\,\tfrac{v^2-p^2}{2r},\,\tfrac vr\bigr),&p\le v\le r,\\[3pt]
\bigl(\tfrac r2+\tfrac{p^2}{2r},\,v-\tfrac r2-\tfrac{p^2}{2r},\,1\bigr),&p\le r\le v.
\end{cases}
\]
Below the reserve the challenger cannot bid, so the outcome is that of no entry. If \(r\le p\le v\), the incumbent
never meets the reserve, and the challenger buys alone at \(p\). The last two rows are the integrals of
\Cref{sec:03_auction-terms} with \(\ell\) or \(h\) replaced by \(v\). The pieces agree at \(v=r\) and at \(p=r\), and a
class law averages the kernel.

\paragraph{Exclusion can widen the spread.} For point values with \(\ell<p<r<h\), the low class adds nothing:
\(t_L=t_0=p-p^2/r\) and \(g_L=0\), while \(t_H=r/2+p^2/(2r)\). So \(\DT=r/2-p+3p^2/(2r)\), the spread of the worked
example in \Cref{sec:03_auction-prop1c}. For \(r<p<h\), \(t_0=t_L=g_L=0\), \(t_H=p\), \(g_H=h-p\) and \(\DT=p\). A larger
spread can make full orders pay even against the weak incumbent.

\paragraph{Preparation, sale and competition differ.} Let \(I\) indicate preparation and \(a=\Pr(I=1,\theta\ge p)\).
Entry is \(\Pr(I=1)\). Two admissible bidders meet with probability \((1-p/r)_+a\), because \(R\) is independent of
\((I,\theta)\), and a sale occurs with probability \((1-p/r)_++a-(1-p/r)_+a\). The zero-cost challenger prepares by the
model's convention, so preparation can occur without an admissible bid.

\begin{resultbox}[Reserves \paper{Prop.~A.10}]
Fix a reserve. Keep the class order, the costs, the slot and a strengthening \(F_1\le F_0\), and compute each class
object with the reserve. Suppose \(k<\rho\kappa\min\{\Delta_0,\Delta_1\}\).
\begin{enumerate}[label=(\alph*),leftmargin=2em]
\item If \(B_0(\pi)<c_H<B_1(M)\), entry is flat: \(\Ent_0=\Ent_1=\rho+(1-\rho)(b+\pi\lambda)\).
\item If \(B_0(\pi)<c_H<B_0(M)\) and \(B_1(M)<c_H\), entry falls strictly:
\(\Ent_0=\rho+(1-\rho)(b+\pi\lambda)>\rho=\Ent_1\).
\end{enumerate}
Each economy has a unique outcome under every pair of tie rules.
\end{resultbox}

\emph{Derivation.} The trading bound of \Cref{sec:05_investor-bounds} needs only \(\DT\ge0\) and entry in \([\rho,1]\),
so it holds with the reserve kernel. Full orders are forced at both strengths, and the posteriors are \((m,\pi,M)\).
Gross profit falls with strength at every belief, because \((\theta-\max\{p,R\})_+\) is nonincreasing in \(R\). In
part (a), \(B_1(\pi)\le B_0(\pi)<c_H<B_1(M)\le B_0(M)\), so only a purchase recruits, at both strengths. In part (b),
the weak strength has the same pattern, and at the strong strength \(B_1(\mu)\le B_1(M)<c_H\) at every posterior. A
purchase has probability \(b+\pi\lambda\) (\Cref{sec:06_main_result-entry}). Every comparison is strict, so no tie rule
matters.

\begin{center}\small
\begin{tabular}{@{}llcccc@{}}\toprule
Reserve & Strength & Preparation & Sale & Two admissible & Proceeds\\\midrule
\(\reg{prim_p}\) & \(r_0\) & \reg{p2_reserve_points_base_r0_entry} & \reg{p2_reserve_points_base_r0_sale} & \reg{p2_reserve_points_base_r0_two} & \reg{p2_reserve_points_base_r0_proceeds}\\
& \(r_1\) & \reg{p2_reserve_points_base_r1_entry} & \reg{p2_reserve_points_base_r1_sale} & \reg{p2_reserve_points_base_r1_two} & \reg{p2_reserve_points_base_r1_proceeds}\\
\(\reg{rsv_points_excl}\) & \(r_0\) & \reg{p2_reserve_points_exclude_r0_entry} & \reg{p2_reserve_points_exclude_r0_sale} & \reg{p2_reserve_points_exclude_r0_two} & \reg{p2_reserve_points_exclude_r0_proceeds}\\
& \(r_1\) & \reg{p2_reserve_points_exclude_r1_entry} & \reg{p2_reserve_points_exclude_r1_sale} & \reg{p2_reserve_points_exclude_r1_two} & \reg{p2_reserve_points_exclude_r1_proceeds}\\
\(\reg{rsv_deter}\) & \(r_0\) & \reg{p2_reserve_points_deter_r0_entry} & \reg{p2_reserve_points_deter_r0_sale} & \reg{p2_reserve_points_deter_r0_two} & \reg{p2_reserve_points_deter_r0_proceeds}\\
& \(r_1\) & \reg{p2_reserve_points_deter_r1_entry} & \reg{p2_reserve_points_deter_r1_sale} & \reg{p2_reserve_points_deter_r1_two} & \reg{p2_reserve_points_deter_r1_proceeds}\\\bottomrule
\end{tabular}
\end{center}
The rows at \(p=\reg{prim_p}\) repeat Proposition~2, with no trade at \(r_0\). The reserve \(\reg{rsv_points_excl}\)
excludes the low class, and part (a) gives flat entry. Entry is exactly flat
because the probability of a purchase, \(b+\pi\lambda\), does not depend on strength. The reserve
\(\reg{rsv_deter}\) gives strict deterrence by part (b); it exceeds \(r_0\), so it also excludes the weak incumbent.

\paragraph{Where strict deterrence holds.} With the low class excluded, \(g_L=0\), so gross profit at the top posterior
is \(Mg_H\). At \(r_1\) it equals the cost when
\[
M\Bigl(h-\frac{r_1}{2}-\frac{p^2}{2r_1}\Bigr)=c_H\iff p=p_H:=\sqrt{2r_1(h-c_H/M)-r_1^2},
\]
and \(B_1(M)<c_H\) above \(p_H\). At \(r_0<p\), \(g_H=h-p\), so \(B_0(M)>c_H\) exactly when \(p<h-c_H/M\). Strict deterrence
therefore holds for \(p_H\approx\reg{p2_reserve_ph}<p<h-c_H/M\approx\reg{p2_reserve_weak_top}\). The other hypotheses of
part (b) hold there: \(B_0(1/2)=(h-p)/2<c_H\), and both spreads rise with \(p\), so the margin of the trading bound,
\(\rho\kappa\min\{\Delta_0,\Delta_1\}-k\), is at least its value \(0.050792\) at \(p_H\). % verify/11_welfare.py
Both endpoints exceed \(r_0\). For \(\ell<p\le r_0\), the top posterior still recruits at both strengths, and entry is
flat. At either endpoint the challenger's tie rule matters. None of this identifies the seller's optimal reserve.
